% Осень26!B16; нумерация внутри столбца без пустых строк.
% Источник: лекция 4, §9 (пространство с мерой)
% Источник: лекция 2, §11 (строка 313)
\definitionticket{13}{Мера, пространство с мерой}
\label{definition:13}

\subsection*{Мера: определение ---}

\begin{definition}

\emph{мера.} Функция \(\mu:\mathcal{P}\to\overline{\mathbb{R}}\), заданная на полукольце («вообще для меры естественно требовать σ-алгебру, но по минимуму подойдёт и полукольцо»), называется \textbf{мерой}, если \(\mu\) --- \textbf{объём} и \(\mu\) \textbf{счётно аддитивна}: \[\forall A,A_1,A_2,\dots\in\mathcal{P}:\qquad A=\bigsqcup_{i=1}^{\infty}A_i\quad\Rightarrow\quad\mu A=\sum_{i=1}^{\infty}\mu A_i.\]

Замечания лектора:

\begin{itemize}
\tightlist
\item
  Оговорка о принадлежности \(A=\bigsqcup A_i\) полукольцу --- как и у конечной аддитивности; в σ-алгебре она автоматична.
\item
  Справа --- \textbf{неотрицательный} ряд, его сумма всегда определена (возможно, \(+\infty\)): «всё максимально просто».
\item
  Нумерация слагаемых не важна: сумма счётного семейства неотрицательных чисел определена безотносительно порядка; можно формулировать и «усиленно» --- для произвольного счётного дизъюнктного семейства, дающего в объединении \(A\); содержательно это то же самое.
\end{itemize}

\end{definition}

\textbf{Пространство с мерой.} Тройка \((X,\mathcal A,\mu)\), где \(X\) --- множество, \(\mathcal A\) --- \(\sigma\)-алгебра подмножеств \(X\), \(\mu\) --- мера на \(\mathcal A\), называется пространством с мерой.
